% !TeX root = ../applied-mathematics1-companion.tex
\chapter{Mechanics: Simple Forces}\label{chap:applied-mechanics}

\section{Newton's laws and inertial frames}
The models in this chapter treat bodies as point particles of constant mass.
Position, velocity, acceleration and force are vectors; in one-dimensional
motion we use signed scalar components along a chosen positive axis.

\begin{definition}[Inertial frame]
An inertial frame is a reference frame in which a force-free particle moves
with constant velocity. Newton's laws below are used in such a frame.
\end{definition}
\begin{enumerate}
\item \textbf{First law.} An isolated particle moves in a straight line at
constant velocity. Zero velocity is just one possibility.
\item \textbf{Second law.} For fixed mass $m$, the net force is
$\mathbf F_{\mathrm{net}}=m\mathbf a=m\ddot{\mathbf r}$.
\item \textbf{Third law.} The force on particle 1 due to particle 2 equals
minus the force on particle 2 due to particle 1:
$\mathbf F_{1\leftarrow2}=-\mathbf F_{2\leftarrow1}$.
The strong form used for central interactions also places them along the
line joining the particles.
\end{enumerate}
Third-law force pairs act on different bodies: they do not cancel on a
free-body diagram of just one body. Multiple forces acting on the same body
must be summed before applying the second law.
In an accelerating frame, extra inertial-force terms are needed; one cannot
simply reuse the same net-force equation unchanged.

\section{Balancing forces}
Near Earth's surface, weight has magnitude $mg$, directed downwards, with
$g\approx9.8\ \mathrm{m\,s^{-2}}$. Mass is measured in kilograms, force in
newtons, and $1\ \mathrm N=1\ \mathrm{kg\,m\,s^{-2}}$.
Weight depends on the gravitational field; mass does not disappear when
an object is weightless.

A taut light string exerts a tension along itself, pulling on the body to
which it is attached. A smooth contact surface exerts a normal reaction
perpendicular to the surface. For a light string over an ideal smooth
massless pulley, the tension magnitude is the same on both sides.
A string cannot push; if a calculation requires negative tension, the taut
string model must be reconsidered.

Static equilibrium requires zero acceleration and zero initial velocity.
Zero net force alone gives constant velocity, which need not be rest.
To balance forces on a particle, resolve the vector equation
$\sum\mathbf F=\mathbf0$ into convenient perpendicular directions.

\studyguide{Draw a separate free-body diagram for each body or junction.
Choose axes and mark the angles from those axes. Resolve every force,
then apply $\sum F=ma$ along each axis. Do not add a body's own third-law
reaction force to its own diagram.}

\begin{example}[Three supporting cables: source Example 3.1]
A mass $m$ hangs from a vertical cable, attached at a light junction to two
inclined cables making angles $\theta_1,\theta_2$ above the horizontal.
Let the inclined tensions be $T_1,T_2$ and the vertical one $T_3$.
% !TeX root = ../applied-mathematics1-companion.tex
\begin{center}
\begin{tikzpicture}[>=Stealth,scale=.9]
\draw[muted,dashed] (-2.4,0)--(2.6,0);
\draw[->,very thick,ink] (0,0)--(-2,1.5) node[above] {$T_1$};
\draw[->,very thick,ink] (0,0)--(2.2,1) node[above] {$T_2$};
\draw[->,very thick,accent] (0,0)--(0,-1.6) node[right] {$T_3=mg$};
\draw (-.7,0) arc[start angle=180,end angle=143.13,radius=.7];
\draw (.8,0) arc[start angle=0,end angle=24.44,radius=.8];
\node at (-1,.29) {$\theta_1$};
\node at (1.1,.22) {$\theta_2$};
\fill (0,0) circle (2pt);
\end{tikzpicture}

{\small Free-body diagram of the light junction; angles are from horizontal.}
\end{center}

For the mass, equilibrium gives $T_3=mg$.
For the junction, resolve horizontally and vertically:
\[
T_1\cos\theta_1=T_2\cos\theta_2,\qquad
T_1\sin\theta_1+T_2\sin\theta_2=mg.
\]
Eliminating one tension and using the sine addition identity gives
\[
T_1=\frac{mg\cos\theta_2}{\sin(\theta_1+\theta_2)},\qquad
T_2=\frac{mg\cos\theta_1}{\sin(\theta_1+\theta_2)},\qquad T_3=mg.
\]
For equal angles $\theta$, both inclined tensions equal $mg/(2\sin\theta)$.
Nearly horizontal cables need very large tension, consistent with the
small vertical components of their forces.
\end{example}

\begin{example}[Blocks on smooth inclines: source Example 3.2]
Two masses $m_1=1\ \mathrm{kg}$ and $m_2=\sqrt2\ \mathrm{kg}$ are joined
by a taut light string over an ideal pulley. The smooth inclines have
angles $\pi/4$ and $\theta$ above horizontal.
For equilibrium, resolve along each plane:
\[
T=m_1g\sin(\pi/4)=\frac g{\sqrt2},\qquad
T=m_2g\sin\theta=\sqrt2g\sin\theta.
\]
Thus $\sin\theta=1/2$; with $0<\theta<\pi/2$, $\theta=\pi/6$.
Perpendicular resolution gives
\[
R_1=m_1g\cos(\pi/4)=\frac g{\sqrt2},\qquad
R_2=m_2g\cos(\pi/6)=g\sqrt{\frac32}.
\]
With masses substituted numerically in kilograms, these force values are
in newtons. The reaction is not generally $mg$: only the perpendicular
component of weight is balanced by it.
\end{example}

\section{Hooke's law and connected springs}
A spring of natural length $\ell$ and stiffness $k>0$ exerts a restoring
force proportional to signed extension $e$:
\[
\text{spring tension}=ke,\qquad\text{force on a displaced end opposes }e.
\]
Stiffness has units $\mathrm{N/m}$. Positive extension means stretching;
negative extension means compression, when a spring pushes outwards.
Hooke's law is an approximation within the spring's elastic range.
An elastic string differs: it supplies tension when extended but goes
slack when shorter than its natural length.

\begin{example}[Three vertical springs: source Example 3.3]
Three identical light springs of natural length $\ell$ and stiffness $k$
connect a ceiling to a floor $3\ell$ below it. A mass $m$ is attached at
each of the two intermediate junctions.
Measure distance downwards from the ceiling. Let the mass positions be
$\ell+x_1$ and $2\ell+x_2$. The signed spring extensions are
\[
e_1=x_1,\qquad e_2=x_2-x_1,\qquad e_3=-x_2.
\]
For the upper mass, downward forces minus upward forces give
\[
mg+k(x_2-x_1)-kx_1=0.
\]
For the lower mass, the middle spring force is upward with signed magnitude
$k(x_2-x_1)$, and the bottom spring contributes downward signed force
$-kx_2$. Thus
\[
mg-k(x_2-x_1)-kx_2=0.
\]
The two equations become
\[
2x_1-x_2=\frac{mg}{k},\qquad2x_2-x_1=\frac{mg}{k}.
\]
Subtracting gives $x_1=x_2$, and hence both equal $mg/k$.
The mass positions below the ceiling are
\[
\ell+\frac{mg}{k},\qquad2\ell+\frac{mg}{k}.
\]
The spring lengths are $\ell+mg/k$, $\ell$, and $\ell-mg/k$.
The bottom spring is compressed and pushes upwards on the lower mass;
the middle spring has zero extension. The ideal geometry requires
$mg/k<\ell$, and physical validity also requires staying in the elastic range.
\end{example}

\section{Equations of one-dimensional motion}
Choose a signed coordinate $x(t)$ and write
\[
v=\dot x,\qquad a=\ddot x=\frac{dv}{dt},\qquad
m\ddot x=F(t,x,\dot x).
\]
Where velocity can be considered locally as a function of position, the
chain rule gives a third useful acceleration form:
\[
a=\frac{dv}{dt}=\frac{dv}{dx}\frac{dx}{dt}=v\frac{dv}{dx}.
\]
Use $dv/dt$ for a velocity--time relation and $v\,dv/dx$ for a
velocity--position relation. Turning points may require treating separate
branches of $v(x)$; one cannot divide by $v$ indiscriminately at $v=0$.

\section{Constant force and time-dependent force}
\subsection{Constant acceleration}
For $\dot v=a$ constant, with $x(0)=x_0$ and $v(0)=u$, integrate twice:
\[
v(t)=u+at,\qquad x(t)=x_0+ut+\frac12at^2.
\]
Alternatively, integrating $v\,dv/dx=a$ gives
\[
v^2=u^2+2a(x-x_0).
\]
These formulae require constant acceleration and consistent signed directions.

\begin{example}[A sinusoidal force: source Example 3.4]
Suppose $F(t)=mA\sin(\omega t)$, with $\omega>0$, $x(0)=0$ and
$v(0)=v_0$. Here $A$ is an acceleration amplitude, not a force.
The equation $\ddot x=A\sin(\omega t)$ integrates to
\begin{align*}
v(t)&=v_0+\frac A\omega\bigl(1-\cos(\omega t)\bigr),\\
x(t)&=\left(v_0+\frac A\omega\right)t
-\frac A{\omega^2}\sin(\omega t).
\end{align*}
The sine term stays bounded, while the first term is linear in $t$.
Consequently $x(t)\to+\infty$ exactly when $v_0>-A/\omega$;
it tends to $-\infty$ for $v_0<-A/\omega$, and remains bounded at equality.
Divergence of position to $+\infty$ does not necessarily mean the velocity
is positive at every sufficiently late instant: bounded oscillations can
remain on top of the drift.
\end{example}

\section{Quadratic resistance: direction matters}
A drag force must oppose velocity. With constant $k>0$, quadratic drag per
unit mass is represented in one dimension by $-kv|v|$.
Choose upward displacement $y$ and upward signed velocity $v=\dot y$.
With gravity, a single signed equation is
\[
\dot v=-g-kv|v|.
\]
During ascent ($v>0$), it becomes $\dot v=-g-kv^2$.
During descent ($v<0$), it becomes $\dot v=-g+kv^2$.
Using $-kv^2$ for both upward and downward velocities would make drag point
downward in both cases, which is physically wrong.
Here $k$ has units inverse length.

\begin{example}[Maximum height: source Example 3.5]
Launch from $y=0$ with $v(0)=u>0$. At the highest point $y=h$, $v=0$.
For the ascending part use $v\,dv/dy=-g-kv^2$, so
\begin{align*}
h&=-\int_u^0\frac{v\,dv}{g+kv^2}\\
&=\frac1{2k}\log\left(\frac{g+ku^2}{g}\right)
=\frac1{2k}\log\left(1+\frac{ku^2}{g}\right).
\end{align*}
As $k\to0$, $\log(1+z)\sim z$ recovers the drag-free height $u^2/(2g)$.
Only ascent was used; the drag sign changes after the turning point.
\end{example}

\begin{example}[Downward terminal speed: source Example 3.6]
Now choose displacement positive \emph{downwards} and let $v(0)=v_0\geq0$.
Gravity is positive and drag negative, so
\[
\dot v=g-kv^2=k(V^2-v^2),\qquad V=\sqrt{\frac gk}.
\]
The equilibrium $v=V$ is the terminal speed. If $v_0=V$, the speed stays
constant. Otherwise, separation gives
\[
\int_{v_0}^{v(t)}\frac{dw}{V^2-w^2}=kt,
\qquad
\frac1{V^2-w^2}=\frac1{2V}\left(\frac1{V-w}+\frac1{V+w}\right).
\]
Hence
\[
\log\left|\frac{V+v}{V-v}\frac{V-v_0}{V+v_0}\right|=2Vkt.
\]
A solution cannot cross the equilibrium $V$ by uniqueness, so $V-v$ keeps
the sign of $V-v_0$. Rearranging gives the expression valid for all
$v_0\geq0$:
\begin{equation}\label{eq:applied-quadratic-drag}
v(t)=V\frac{(V+v_0)e^{2Vkt}-(V-v_0)}
              {(V+v_0)e^{2Vkt}+(V-v_0)}.
\end{equation}
It also includes $v_0=V$ by direct substitution. For $v_0<V$ the speed
increases towards $V$; for $v_0>V$ it decreases towards $V$.
Equivalently, when $v_0<V$,
\[
v(t)=V\tanh\left(Vkt+\operatorname{artanh}(v_0/V)\right).
\]
For $v_0>V$, the corresponding hyperbolic form uses $\coth$, not a real
$\operatorname{artanh}(v_0/V)$. The exponential fraction avoids this case split.
\end{example}

\section{Linear resistance and terminal velocity}
For constant applied force per unit mass $a_0$ and linear drag coefficient
$\lambda>0$, the signed equation is
\[
\dot v=a_0-\lambda v.
\]
The drag term $-\lambda v$ automatically reverses with the velocity.
The terminal velocity is $v_\infty=a_0/\lambda$, which may be positive,
negative or zero. The coefficient $\lambda$ has units inverse time.
An integrating factor gives, for every real $v_0$,
\begin{align*}
v(t)&=v_\infty+(v_0-v_\infty)e^{-\lambda t},\\
x(t)&=x_0+v_\infty t+
\frac{v_0-v_\infty}{\lambda}(1-e^{-\lambda t}).
\end{align*}
The difference from terminal velocity preserves its sign and decays
exponentially. Starting exactly at terminal velocity produces uniform motion.

\begin{example}[A bullet in a block: source Example 3.8]
A particle enters a resistive block at $x=0$ with speed $u>0$. Resistance
per unit mass is $kv$, with $k>0$, and gravity is neglected.
Then $\dot v=-kv$ and
\[
v(t)=ue^{-kt},\qquad x(t)=\frac uk(1-e^{-kt}).
\]
The limiting penetration is $u/k$. It is approached as $t\to\infty$;
in this ideal model, the particle does not stop exactly in finite time.
A block of width at least $u/k$ prevents finite-time escape through the
far side. To reach half the limiting distance,
\[
\frac u{2k}=\frac uk(1-e^{-kT})
\quad\Longrightarrow\quad T=\frac{\log2}{k}.
\]

Now two particles enter simultaneously from opposite ends, initially
separated by $L>0$, with inward speeds $u_1,u_2>0$.
Their distances into the block at time $t$ are
\[
s_1(t)=\frac{u_1}{k}(1-e^{-kt}),\qquad
s_2(t)=\frac{u_2}{k}(1-e^{-kt}).
\]
Finite-time collision requires $s_1+s_2=L$. Since $1-e^{-kt}<1$ for every
finite $t\geq0$, the necessary and sufficient condition is
\[
\boxed{u_1+u_2>kL}.
\]
When $u_1+u_2=kL$, the separation tends to zero but is positive at every
finite time; that is a limiting meeting, not a finite-time collision.
This corrects the non-strict inequality in both supplied versions.
For a finite collision,
\[
T=\frac1k\log\frac{u_1+u_2}{u_1+u_2-kL},\qquad
s_1(T)=\frac{u_1L}{u_1+u_2}.
\]
The second expression gives the collision position measured from the
first particle's entry end. As $u_1+u_2$ decreases to $kL$ from above,
the collision time diverges, confirming the strict threshold.
\end{example}

\studyguide{For motion: choose the positive direction first; draw forces;
write the signed net force; choose $dv/dt$ or $v\,dv/dx$ according to the
relation wanted; include initial conditions and equilibrium velocities;
check the result's units, signs and limiting behaviour.}
